\section{Experimental Design}

We organize the study around four questions.

\begin{description}
    \item[RQ1: Synthesis.] Can the system derive applicable, executable
    quantitative invariants from public task contracts without encoding hidden
    answers?
    \item[RQ2: Search.] Does an invariant signature help the Evolver select or
    synthesize more useful harness components than matched generic trajectory
    evidence alone?
    \item[RQ3: Improvement.] Do the resulting candidates improve official task
    outcomes, and which gains repeat or transfer under matched invariant
    failures?
    \item[RQ4: Efficiency.] Does invariant-guided search reach useful
    candidates with fewer Worker/verifier calls or lower total model cost?
\end{description}

\subsection{Benchmarks and Evidence Splits}

The planned evaluation uses QuantCodeEval and QFBench. Any task whose feedback
affects candidate generation or selection is optimization/development evidence.
The final task panel is evaluated only after the harness and method are frozen,
and its results do not return to the Evolver.

\subsection{Baselines}

The planned comparison includes:

\begin{itemize}
    \item the frozen minimal Worker harness $H_0$;
    \item prompt-only evolution;
    \item matched generic full-harness evolution using trajectories and official
    outcomes but no invariant signature;
    \item a faithful AHE- or Meta-Harness-style quant adaptation under matched
    model, evidence, mutation surface, and budget;
    \item the full invariant-guided method.
\end{itemize}

Where practical, we also compare a fixed-harness research-state baseline and a
compute-matched repeated-Worker baseline. We report model calls, tokens, audit
calls, official verifier calls, runtime, and cost rather than matching only one
resource axis.

\subsection{Ordered Mechanism Route}

The current prospective order is AP-2M (feedback-driven autonomous probe), QI-1
(task-conditioned invariant synthesis), QI-2 (invariant-guided component
search), and AP-3 (bootstrap from the minimal harness using invariant evidence
only if QI-1/QI-2 help). These are mechanism canaries; the paper-level result
requires repeated campaigns and frozen evaluation.
